% !Mode:: "TeX:UTF-8"
% !TEX program  = xelatex
\section{引言}
随着大语言模型的广泛应用，基于 LoRA 的技术，被广泛用于支持多任务与个性化推理。然而，在实际部署中，一个推理系统往往需要同时服务大量 LoRA 适配器，每个请求依赖于不同的适配器，这对系统的资源管理能力提出了新的挑战。

在多 LoRA 推理场景中，通常仅有一个基础模型常驻 GPU，而大量 LoRA 适配器存储在 CPU 或外部存储中。由于 GPU 显存有限，系统无法同时加载所有适配器，因此需要在运行过程中动态决定哪些适配器驻留在 GPU 上，哪些适配器需要被替换或重新加载。频繁的适配器加载与切换会带来显著的开销，从而影响系统的整体性能。

现有方法大多采用被动的适配器管理策略，例如基于请求触发的加载机制或简单的缓存策略（如 FIFO 或 LRU）。然而，这类方法缺乏对未来请求的预判能力，容易导致频繁的缓存失效与适配器切换。为了解决上述问题，本文提出一种基于预测模型的适配器管理方法。该方法通过分析历史请求模式与系统状态，对未来适配器的使用情况进行估计，并据此为每个适配器分配一个效用评分，从而指导适配器的驻留与淘汰决策。在显存受限的条件下，该方法能够优先保留高价值适配器，从而减少加载开销，提高系统性能。
